\subsection*{CU-01: Iniciar sesión}
\addcontentsline{toc}{subsection}{CU-01: Iniciar sesión}
\label{cu-01}

\begin{fichacu}{CU-01}{Iniciar sesión}
  \crudFila{Responsable}{José Eduardo Prior Hernández}
  \crudFila{Actor principal}{Jugador}
  \crudFila{Actores de sistema}{Servidor de partidas, Servicio de correo, Proveedor del segundo factor, Base de datos}
  \crudFila{Prototipos}{6c, 7c, 7e}
  \crudFila{Objetivo}{Permitir al Jugador acreditar su identidad ante el SJM para obtener una sesión activa en su dispositivo, con el objetivo de acceder al menú principal y a las funciones que exigen una cuenta: jugar partidas en línea, consultar sus \textit{EstadisticaModo}, administrar sus objetos cosméticos y participar en el chat. El caso concentra además las barreras de acceso del sistema: el bloqueo temporal por intentos fallidos, el estado de la cuenta, las sanciones y el segundo factor de autenticación.}
  \crudFila{Disparador}{El Jugador selecciona la opción ``Iniciar sesión'' en la \texttt{GUI\_\allowbreak{}MainScreen}, o bien el cliente del SJM se inicia sin una sesión válida almacenada.}
  \textbf{Precondiciones} & \cuCelda{\begin{cureglas}
      \item \textbf{\mbox{PRE-01}} El Jugador tiene una cuenta registrada en el SJM; es decir, existe un registro de \textit{Usuario} con \texttt{tipo\_\allowbreak{}cuenta} en \texttt{REGISTRADA} asociado a él.
      \item \textbf{\mbox{PRE-02}} El cliente del SJM está instalado y en ejecución en el dispositivo del Jugador.
      \item \textbf{\mbox{PRE-03}} El Servidor de partidas está en ejecución y acepta conexiones TCP en el extremo de red configurado.
      \item \textbf{\mbox{PRE-04}} El Servidor de partidas tiene una conexión disponible con la Base de datos.
      \item \textbf{\mbox{PRE-05}} El cliente no tiene una sesión abierta; si conserva un token de sesión previo, este ya expiró o fue invalidado.
    \end{cureglas}} \\ \hline
  \textbf{Flujo normal} & \cuCelda{\begin{cuenum}[start=1]
      \item El SJM despliega la \texttt{GUI\_\allowbreak{}Login}, solicitando el identificador (\texttt{nickname} o \texttt{correo}) y la contraseña del \textit{Usuario}, y mostrando las opciones ``Iniciar sesión'', ``Crear cuenta'', ``Jugar como invitado'', ``Olvidé mi contraseña'', el selector de idioma y ``Salir''. \textit{(Ver \mbox{FA-01}, \mbox{FA-02}, \mbox{FA-03}, \mbox{FA-04})}
      \item El Jugador captura su identificador y su contraseña, y da click en ``Iniciar sesión''.
      \item El SJM valida en el cliente que ambos campos estén completos y que el identificador cumpla el formato de \texttt{nickname} o de \texttt{correo}. \textit{(Ver \mbox{FA-05})}
      \item El SJM abre la conexión TCP con el Servidor de partidas y negocia la versión del protocolo. \textit{(Ver \mbox{EX-01}, \mbox{EX-02})}
      \item El SJM envía al Servidor de partidas el mensaje \texttt{LOGIN\_\allowbreak{}REQUEST} con el identificador, la contraseña, la huella del dispositivo y la \texttt{version\_\allowbreak{}cliente}. \textit{(Ver \mbox{EX-03}, \mbox{EX-07}, \mbox{EX-08})}
      \item El Servidor de partidas recupera de la Base de datos el registro de \textit{Usuario} cuyo \texttt{nickname} o \texttt{correo} coincida con el identificador recibido. \textit{(Ver \mbox{EX-04})}
      \item El Servidor de partidas verifica que el \textit{Usuario} exista y que su \texttt{tipo\_\allowbreak{}cuenta} sea \texttt{REGISTRADA}. \textit{(Ver \mbox{FA-06})}
      \item El Servidor de partidas verifica que el \textit{Usuario} no tenga un bloqueo temporal vigente, comparando \texttt{bloqueada\_\allowbreak{}hasta} con la fecha y hora actuales. \textit{(Ver \mbox{FA-07})}
    \end{cuenum}} \\
   & \cuCelda{\begin{cuenum}[start=9]
      \item El Servidor de partidas deriva el resumen criptográfico de la contraseña recibida utilizando la \texttt{contrasena\_\allowbreak{}sal} almacenada del \textit{Usuario} y lo compara contra \texttt{contrasena\_\allowbreak{}hash}. \textit{(Ver \mbox{FA-08})}
      \item El Servidor de partidas restablece \texttt{intentos\_\allowbreak{}fallidos} a cero y deja \texttt{bloqueada\_\allowbreak{}hasta} en nulo.
      \item El Servidor de partidas verifica que el \texttt{estado\_\allowbreak{}cuenta} del \textit{Usuario} sea \texttt{ACTIVA}. \textit{(Ver \mbox{FA-09}, \mbox{FA-10}, \mbox{FA-19})}
      \item El Servidor de partidas consulta si existe una \textit{Sancion} del \textit{Usuario} con \texttt{ambito} igual a \texttt{CUENTA}, \texttt{activa} verdadero y \texttt{fecha\_\allowbreak{}fin} nula o posterior a la fecha actual. \textit{(Ver \mbox{FA-10})}
      \item El Servidor de partidas consulta el valor de \texttt{doble\_\allowbreak{}factor\_\allowbreak{}habilitado} del \textit{Usuario} y confirma que está desactivado. \textit{(Ver \mbox{FA-11})}
      \item El Servidor de partidas genera un token de sesión y registra una nueva \textit{Sesion} asociada al \textit{Usuario}, con su \texttt{fecha\_\allowbreak{}inicio}, su \texttt{fecha\_\allowbreak{}expiracion}, la \texttt{direccion\_\allowbreak{}ip} de origen, la \texttt{huella\_\allowbreak{}dispositivo}, el \texttt{nombre\_\allowbreak{}dispositivo} y la \texttt{version\_\allowbreak{}cliente}. \textit{(Ver \mbox{EX-05})}
      \item El Servidor de partidas registra el acceso exitoso en la \textit{BitacoraAcceso} con el resultado \texttt{EXITO}, en la misma transacción que la \textit{Sesion}.
      \item El Servidor de partidas confirma la transacción. \textit{(Ver \mbox{EX-05})}
    \end{cuenum}} \\
   & \cuCelda{\begin{cuenum}[start=17]
      \item El Servidor de partidas responde con el mensaje \texttt{LOGIN\_\allowbreak{}OK}, que contiene el token de sesión, el perfil del \textit{Usuario} (\texttt{nickname}, \texttt{id\_\allowbreak{}icono}, \texttt{idioma\_\allowbreak{}preferido}), su \textit{Equipamiento}, sus \textit{EstadisticaModo} y, si la tiene, la \textit{Participacion} sin terminar.
      \item El SJM aplica el \texttt{idioma\_\allowbreak{}preferido} recibido a los textos de la interfaz. \textit{(Ver \mbox{FA-17}, \mbox{FA-18})}
      \item El SJM despliega la \texttt{GUI\_\allowbreak{}MainMenu} con el \texttt{nickname}, el peón equipado y el elo del modo clásico 9×9.
      \item Fin del caso de uso.
    \end{cuenum}} \\ \hline
  \textbf{Flujos alternos} & \cuCelda{\textbf{\mbox{FA-01}: Salir del juego}\par
    \begin{cuenum}[start=1]
      \item El Jugador selecciona ``Salir''.
      \item El SJM cierra la conexión TCP, si estuviera abierta, y la aplicación.
      \item Fin del caso de uso.
    \end{cuenum}} \\
   & \cuCelda{\textbf{\mbox{FA-02}: Crear una cuenta o jugar como invitado}\par
    \begin{cuenum}[start=1]
      \item El Jugador selecciona ``Crear cuenta'' o ``Jugar como invitado''.
      \item El flujo continúa en el \mbox{CU-02} Registrar cuenta o en el \mbox{CU-06} Jugar como invitado, respectivamente.
      \item Fin del caso de uso.
    \end{cuenum}} \\
   & \cuCelda{\textbf{\mbox{FA-03}: Recuperar la contraseña}\par
    \begin{cuenum}[start=1]
      \item El Jugador selecciona ``Olvidé mi contraseña''.
      \item El flujo continúa en el \mbox{CU-03} Recuperar contraseña.
      \item Fin del caso de uso.
    \end{cuenum}} \\
   & \cuCelda{\textbf{\mbox{FA-04}: Cambiar el idioma de la interfaz}\par
    \begin{cuenum}[start=1]
      \item El Jugador selecciona un idioma distinto en el selector de idioma.
      \item El SJM carga el catálogo de recursos correspondiente y vuelve a dibujar los textos de la \texttt{GUI\_\allowbreak{}Login}, descartando la contraseña capturada.
      \item El SJM guarda el idioma elegido en la configuración local del cliente y no en la base de datos: esta pantalla ocurre antes de saber quién es el \textit{Usuario}, así que no puede depender de su \texttt{idioma\_\allowbreak{}preferido}.
      \item El flujo continúa en el paso 2 del FN.
    \end{cuenum}} \\
   & \cuCelda{\textbf{\mbox{FA-05}: Campos incompletos o con formato inválido}\par
    \begin{cuenum}[start=1]
      \item El SJM muestra la \texttt{GUI\_\allowbreak{}MessageWarning} indicando qué campos están incompletos o mal formados, y los resalta en el formulario.
      \item El flujo continúa en el paso 2 del FN. La validación es local y no llega al Servidor de partidas, por lo que no incrementa \texttt{intentos\_\allowbreak{}fallidos} (\mbox{RN-13}).
    \end{cuenum}} \\
   & \cuCelda{\textbf{\mbox{FA-06}: El identificador no corresponde a ninguna cuenta registrada}\par
    \begin{cuenum}[start=1]
      \item El Servidor de partidas ejecuta de todos modos una derivación de resumen criptográfico sobre una sal ficticia, de forma que el tiempo de respuesta no permita distinguir una cuenta inexistente de una contraseña equivocada (\mbox{RN-01}).
      \item El Servidor de partidas registra el intento fallido en la \textit{BitacoraAcceso}, con el \texttt{identificador\_\allowbreak{}capturado}, la \texttt{direccion\_\allowbreak{}ip} de origen y el resultado \texttt{USUARIO\_\allowbreak{}INEXISTENTE}, dejando el \texttt{id\_\allowbreak{}usuario} en nulo.
      \item El Servidor de partidas responde con \texttt{LOGIN\_\allowbreak{}FAILED} y el motivo genérico \texttt{CREDENCIALES\_\allowbreak{}INVALIDAS}.
      \item El SJM muestra la \texttt{GUI\_\allowbreak{}MessageWarning} con el mensaje ``El usuario o la contraseña son incorrectos''.
      \item El flujo continúa en el paso 2 del FN.
    \end{cuenum}} \\
   & \cuCelda{\textbf{\mbox{FA-07}: La cuenta tiene un bloqueo temporal vigente}\par
    \begin{cuenum}[start=1]
      \item El Servidor de partidas calcula los minutos restantes entre la fecha y hora actuales y \texttt{bloqueada\_\allowbreak{}hasta}.
      \item El Servidor de partidas registra el intento en la \textit{BitacoraAcceso} con el resultado \texttt{BLOQUEO\_\allowbreak{}VIGENTE}, sin modificar \texttt{intentos\_\allowbreak{}fallidos} ni prolongar el bloqueo, de modo que reintentar no alargue la espera.
      \item El Servidor de partidas responde con \texttt{LOGIN\_\allowbreak{}FAILED}, el motivo \texttt{CUENTA\_\allowbreak{}BLOQUEADA\_\allowbreak{}TEMPORALMENTE} y los minutos restantes.
      \item El SJM muestra la \texttt{GUI\_\allowbreak{}MessageWarning} indicando que la cuenta está bloqueada temporalmente y cuántos minutos faltan.
      \item Fin del caso de uso.
    \end{cuenum}} \\
   & \cuCelda{\textbf{\mbox{FA-08}: La contraseña es incorrecta}\par
    \begin{cuenum}[start=1]
      \item El Servidor de partidas incrementa en uno \texttt{intentos\_\allowbreak{}fallidos} del \textit{Usuario} y actualiza \texttt{fecha\_\allowbreak{}ultimo\_\allowbreak{}intento\_\allowbreak{}fallido} con la fecha y hora actuales.
      \item El Servidor de partidas evalúa \texttt{intentos\_\allowbreak{}fallidos} contra la regla \mbox{RN-03} y, si alcanza un umbral, calcula \texttt{bloqueada\_\allowbreak{}hasta} sumando a la fecha actual la duración correspondiente.
      \item El Servidor de partidas guarda \texttt{intentos\_\allowbreak{}fallidos}, \texttt{fecha\_\allowbreak{}ultimo\_\allowbreak{}intento\_\allowbreak{}fallido} y \texttt{bloqueada\_\allowbreak{}hasta} del \textit{Usuario}. \textit{(Ver \mbox{EX-04})}
      \item El Servidor de partidas registra el intento en la \textit{BitacoraAcceso} con el resultado \texttt{CONTRASENA\_\allowbreak{}INCORRECTA}, sin incluir en ningún caso la contraseña recibida (\mbox{RN-02}).
      \item El Servidor de partidas responde con \texttt{LOGIN\_\allowbreak{}FAILED} y el motivo genérico \texttt{CREDENCIALES\_\allowbreak{}INVALIDAS}; si la cuenta quedó bloqueada, incluye además los minutos de bloqueo.
      \item El SJM muestra la \texttt{GUI\_\allowbreak{}MessageWarning} con el mensaje correspondiente.
      \item Si la cuenta quedó bloqueada, fin del caso de uso. En caso contrario, el flujo continúa en el paso 2 del FN.
    \end{cuenum}} \\
   & \cuCelda{\textbf{\mbox{FA-09}: La cuenta está pendiente de verificación}\par
    \begin{cuenum}[start=1]
      \item El Servidor de partidas registra el intento en la \textit{BitacoraAcceso} con el resultado \texttt{CUENTA\_\allowbreak{}PENDIENTE}.
      \item El Servidor de partidas responde con \texttt{LOGIN\_\allowbreak{}FAILED} y el motivo \texttt{CUENTA\_\allowbreak{}PENDIENTE\_\allowbreak{}VERIFICACION}.
      \item El SJM despliega la \texttt{GUI\_\allowbreak{}PendingVerification}, mostrando el \texttt{correo} registrado de forma parcialmente oculta y las opciones ``Reenviar correo de verificación'' y ``Aceptar''.
      \item El Jugador selecciona ``Reenviar correo de verificación'' y el SJM envía el mensaje \texttt{RESEND\_\allowbreak{}VERIFICATION\_\allowbreak{}REQUEST}. Si selecciona ``Aceptar'', el flujo continúa en el paso 2 del FN.
      \item El Servidor de partidas verifica que hayan transcurrido al menos cinco minutos desde \texttt{fecha\_\allowbreak{}ultimo\_\allowbreak{}envio\_\allowbreak{}verificacion}, o que este valor sea nulo (\mbox{RN-10}). \textit{(Ver \mbox{FA-12})}
      \item El Servidor de partidas invalida los \textit{TokenVerificacionCorreo} pendientes del \textit{Usuario}, genera uno nuevo con su fecha de expiración y lo guarda. \textit{(Ver \mbox{EX-04})}
      \item El Servidor de partidas solicita al Servicio de correo el envío del mensaje de verificación en el \texttt{idioma\_\allowbreak{}preferido} del \textit{Usuario}. \textit{(Ver \mbox{EX-06})}
      \item El Servidor de partidas actualiza \texttt{fecha\_\allowbreak{}ultimo\_\allowbreak{}envio\_\allowbreak{}verificacion} del \textit{Usuario}.
      \item El Servidor de partidas responde con \texttt{RESEND\_\allowbreak{}VERIFICATION\_\allowbreak{}OK}, que el SJM informa mediante la \texttt{GUI\_\allowbreak{}MessageSuccess}.
      \item Fin del caso de uso.
    \end{cuenum}} \\
   & \cuCelda{\textbf{\mbox{FA-10}: La cuenta tiene una sanción de ámbito de cuenta}\par
    \begin{cuenum}[start=1]
      \item El Servidor de partidas recupera de la \textit{Sancion} vigente el \texttt{tipo}, el \texttt{motivo} y la \texttt{fecha\_\allowbreak{}fin}.
      \item El Servidor de partidas registra el intento en la \textit{BitacoraAcceso} con el resultado \texttt{CUENTA\_\allowbreak{}SANCIONADA}.
      \item El Servidor de partidas responde con \texttt{LOGIN\_\allowbreak{}FAILED}, el motivo \texttt{CUENTA\_\allowbreak{}SANCIONADA}, los datos de la \textit{Sancion} y un token de apelación de corta vigencia que identifica la sanción sin abrir sesión.
      \item El SJM despliega la \texttt{GUI\_\allowbreak{}BannedAccount} con el motivo de la sanción, su tipo, la fecha en que termina si es temporal y la opción ``Apelar'', que continúa en el \mbox{CU-45} Apelar una sanción.
      \item Fin del caso de uso. Una sanción de ámbito \texttt{CHAT} no llega a este flujo: no impide iniciar sesión (\mbox{D-05}).
    \end{cuenum}} \\
   & \cuCelda{\textbf{\mbox{FA-11}: La cuenta tiene habilitado el segundo factor}\par
    \begin{cuenum}[start=1]
      \item El Servidor de partidas genera un \textit{CodigoSegundoFactor} de un solo uso con su fecha de expiración, lo asocia al \textit{Usuario} y guarda únicamente su resumen criptográfico. \textit{(Ver \mbox{EX-04})}
      \item El Servidor de partidas entrega el código al Proveedor del segundo factor configurado. \textit{(Ver \mbox{EX-06})}
      \item El Servidor de partidas responde con \texttt{SECOND\_\allowbreak{}FACTOR\_\allowbreak{}REQUIRED} y la fecha de expiración del código.
      \item El SJM despliega la \texttt{GUI\_\allowbreak{}SecondFactor}, solicitando el código y mostrando el tiempo restante de vigencia, con las opciones ``Verificar'', ``Reenviar código'' y ``Cancelar''. \textit{(Ver \mbox{FA-13}, \mbox{FA-14})}
      \item El Jugador captura el código y el SJM envía el mensaje \texttt{SECOND\_\allowbreak{}FACTOR\_\allowbreak{}REQUEST}.
      \item El Servidor de partidas verifica que el código no haya expirado y que no haya sido consumido. \textit{(Ver \mbox{FA-15})}
      \item El Servidor de partidas compara el resumen del código recibido contra el almacenado. \textit{(Ver \mbox{FA-16})}
      \item El Servidor de partidas marca el \textit{CodigoSegundoFactor} como consumido.
      \item El flujo continúa en el paso 14 del FN.
    \end{cuenum}} \\
   & \cuCelda{\textbf{\mbox{FA-12}: Se solicita el reenvío antes del tiempo permitido}\par
    \begin{cuenum}[start=1]
      \item El Servidor de partidas responde con \texttt{RESEND\_\allowbreak{}VERIFICATION\_\allowbreak{}THROTTLED} y los segundos que faltan, que el SJM muestra en la \texttt{GUI\_\allowbreak{}MessageWarning}.
      \item El flujo continúa en el paso 4 del \mbox{FA-09}.
    \end{cuenum}} \\
   & \cuCelda{\textbf{\mbox{FA-13}: Cancelar la verificación del segundo factor}\par
    \begin{cuenum}[start=1]
      \item El Jugador selecciona ``Cancelar'' en la \texttt{GUI\_\allowbreak{}SecondFactor} y el SJM envía el mensaje \texttt{ABORT\_\allowbreak{}LOGIN\_\allowbreak{}REQUEST}.
      \item El Servidor de partidas invalida el \textit{CodigoSegundoFactor} pendiente y registra el abandono en la \textit{BitacoraAcceso} con el resultado \texttt{ABANDONADO}.
      \item El flujo continúa en el paso 2 del FN.
    \end{cuenum}} \\
   & \cuCelda{\textbf{\mbox{FA-14}: Reenviar el código del segundo factor}\par
    \begin{cuenum}[start=1]
      \item El Jugador selecciona ``Reenviar código'' y el SJM envía el mensaje \texttt{RESEND\_\allowbreak{}SECOND\_\allowbreak{}FACTOR\_\allowbreak{}REQUEST}.
      \item El Servidor de partidas verifica que hayan transcurrido al menos sesenta segundos desde la \texttt{fecha\_\allowbreak{}generacion} del código vigente. Si no se cumple, responde con los segundos restantes y el flujo continúa en el paso 5 del \mbox{FA-11}.
      \item El Servidor de partidas invalida el código anterior, genera uno nuevo y lo entrega al Proveedor del segundo factor. \textit{(Ver \mbox{EX-06})}
      \item El flujo continúa en el paso 5 del \mbox{FA-11}.
    \end{cuenum}} \\
   & \cuCelda{\textbf{\mbox{FA-15}: El código del segundo factor expiró}\par
    \begin{cuenum}[start=1]
      \item El Servidor de partidas marca el \textit{CodigoSegundoFactor} como invalidado.
      \item El Servidor de partidas responde con \texttt{SECOND\_\allowbreak{}FACTOR\_\allowbreak{}EXPIRED}, que el SJM informa mediante la \texttt{GUI\_\allowbreak{}MessageWarning}.
      \item El flujo continúa en el paso 5 del \mbox{FA-11}.
    \end{cuenum}} \\
   & \cuCelda{\textbf{\mbox{FA-16}: El código del segundo factor es incorrecto}\par
    \begin{cuenum}[start=1]
      \item El Servidor de partidas incrementa en uno el atributo \texttt{intentos} del \textit{CodigoSegundoFactor}.
      \item El Servidor de partidas evalúa \texttt{intentos} contra la regla \mbox{RN-09}. Si alcanza el máximo, invalida el código, cancela el intento de inicio de sesión y lo registra en la \textit{BitacoraAcceso} con el resultado \texttt{SEGUNDO\_\allowbreak{}FACTOR\_\allowbreak{}FALLIDO}.
      \item El Servidor de partidas responde con \texttt{SECOND\_\allowbreak{}FACTOR\_\allowbreak{}INVALID} y el número de intentos restantes, que el SJM muestra en la \texttt{GUI\_\allowbreak{}MessageWarning}.
      \item Si se agotaron los intentos, el flujo continúa en el paso 2 del FN. En caso contrario, continúa en el paso 5 del \mbox{FA-11}.
    \end{cuenum}} \\
   & \cuCelda{\textbf{\mbox{FA-17}: Es el primer acceso de la cuenta}\par
    \begin{cuenum}[start=1]
      \item El SJM detecta en el perfil recibido que \texttt{fecha\_\allowbreak{}configuracion\_\allowbreak{}inicial} es nula.
      \item El flujo continúa en el \mbox{CU-08} Configurar la cuenta la primera vez, que al terminar despliega la \texttt{GUI\_\allowbreak{}MainMenu}.
      \item Fin del caso de uso.
    \end{cuenum}} \\
   & \cuCelda{\textbf{\mbox{FA-18}: El Jugador tiene una partida sin terminar}\par
    \begin{cuenum}[start=1]
      \item El SJM detecta en el mensaje \texttt{LOGIN\_\allowbreak{}OK} una \textit{Participacion} sin terminar en una partida en línea.
      \item El flujo continúa en el \mbox{CU-26} Reconectar a una partida en curso.
      \item Fin del caso de uso.
    \end{cuenum}} \\
   & \cuCelda{\textbf{\mbox{FA-19}: La suspensión temporal de la cuenta ya venció}\par
    \begin{cuenum}[start=1]
      \item El Servidor de partidas encuentra el \texttt{estado\_\allowbreak{}cuenta} en \texttt{SUSPENDIDA} y ninguna \textit{Sancion} de ámbito \texttt{CUENTA} con \texttt{activa} verdadero y \texttt{fecha\_\allowbreak{}fin} posterior a la fecha actual.
      \item El Servidor de partidas devuelve el \texttt{estado\_\allowbreak{}cuenta} a \texttt{ACTIVA} y registra el cambio en la \textit{BitacoraAcceso} con el resultado \texttt{SUSPENSION\_\allowbreak{}VENCIDA} (\mbox{RN-16}). \textit{(Ver \mbox{EX-04})}
      \item El flujo continúa en el paso 13 del FN. Si la sanción sigue vigente, el flujo es el del \mbox{FA-10}.
    \end{cuenum}} \\ \hline
  \textbf{Excepciones} & \cuCelda{\textbf{\mbox{EX-01}: No se puede establecer la conexión con el servidor}\par
    \begin{cuenum}[start=1]
      \item El SJM detecta que la conexión TCP no pudo establecerse dentro del tiempo límite configurado, y registra el fallo en la bitácora local del cliente.
      \item El SJM muestra la \texttt{GUI\_\allowbreak{}MessageError} indicando que no fue posible contactar al servidor, con las opciones ``Reintentar'' y ``Aceptar''.
      \item El Jugador selecciona ``Reintentar'' y el flujo continúa en el paso 4 del FN. Si selecciona ``Aceptar'', continúa en el paso 2 del FN.
    \end{cuenum}} \\
   & \cuCelda{\textbf{\mbox{EX-02}: La versión del protocolo es incompatible}\par
    \begin{cuenum}[start=1]
      \item El Servidor de partidas responde con \texttt{PROTOCOL\_\allowbreak{}VERSION\_\allowbreak{}MISMATCH}, indicando la versión mínima que acepta, y cierra la conexión.
      \item El SJM muestra la \texttt{GUI\_\allowbreak{}MessageError} indicando que la versión instalada del juego es incompatible y que debe actualizarse.
      \item Fin del caso de uso.
    \end{cuenum}} \\
   & \cuCelda{\textbf{\mbox{EX-03}: Se agota el tiempo de espera de la respuesta}\par
    \begin{cuenum}[start=1]
      \item El SJM detecta que transcurrió el tiempo límite de respuesta sin haber recibido un mensaje completo del Servidor de partidas.
      \item El SJM cierra la conexión TCP y descarta el intercambio en curso, sin suponer que la solicitud fracasó: el servidor pudo haberla procesado.
      \item El SJM registra el fallo en la bitácora local del cliente y lo informa mediante la \texttt{GUI\_\allowbreak{}MessageError}.
      \item El flujo continúa en el paso 2 del FN.
    \end{cuenum}} \\
   & \cuCelda{\textbf{\mbox{EX-04}: Fallo al consultar o actualizar la base de datos}\par
    \begin{cuenum}[start=1]
      \item El Servidor de partidas revierte la transacción en curso, si la hubiera.
      \item El Servidor de partidas registra la excepción en la bitácora del servidor con un identificador de incidencia, sin incluir la contraseña ni el código recibidos (\mbox{RN-02}).
      \item El Servidor de partidas responde con \texttt{SERVER\_\allowbreak{}ERROR} y el identificador de incidencia, y cierra la conexión.
      \item El SJM muestra la \texttt{GUI\_\allowbreak{}MessageError} con el identificador de incidencia, para que el Jugador pueda citarlo al reportar el fallo.
      \item Fin del caso de uso.
    \end{cuenum}} \\
   & \cuCelda{\textbf{\mbox{EX-05}: Fallo al registrar la sesión}\par
    \begin{cuenum}[start=1]
      \item El Servidor de partidas revierte la transacción que abarca el registro de la \textit{Sesion} y el del acceso exitoso en la \textit{BitacoraAcceso}, de modo que no quede un acceso registrado sin sesión.
      \item El Servidor de partidas registra la excepción en la bitácora del servidor y responde con \texttt{SERVER\_\allowbreak{}ERROR}.
      \item El SJM muestra la \texttt{GUI\_\allowbreak{}MessageError} indicando que no fue posible abrir la sesión.
      \item Fin del caso de uso. El Jugador queda sin sesión activa.
    \end{cuenum}} \\
   & \cuCelda{\textbf{\mbox{EX-06}: El servicio externo no está disponible}\par
    \begin{cuenum}[start=1]
      \item El Servidor de partidas agota los reintentos previstos hacia el Servicio de correo o el Proveedor del segundo factor sin obtener confirmación de entrega.
      \item El Servidor de partidas invalida el código o el token que había generado, para que no quede uno vigente que nadie recibió.
      \item El Servidor de partidas registra el fallo del servicio externo en la bitácora del servidor y responde con \texttt{EXTERNAL\_\allowbreak{}SERVICE\_\allowbreak{}UNAVAILABLE}.
      \item El SJM muestra la \texttt{GUI\_\allowbreak{}MessageError} indicando que no fue posible enviar el código o el correo.
      \item Fin del caso de uso.
    \end{cuenum}} \\
   & \cuCelda{\textbf{\mbox{EX-07}: La conexión se interrumpe durante el intercambio}\par
    \begin{cuenum}[start=1]
      \item El SJM detecta el cierre inesperado del socket antes de recibir un mensaje completo.
      \item El SJM descarta el token de sesión que hubiera recibido de forma parcial y registra el fallo en la bitácora local del cliente.
      \item El SJM muestra la \texttt{GUI\_\allowbreak{}MessageError} indicando que se perdió la conexión con el servidor.
      \item El flujo continúa en el paso 2 del FN.
    \end{cuenum}} \\
   & \cuCelda{\textbf{\mbox{EX-08}: El mensaje recibido está malformado}\par
    \begin{cuenum}[start=1]
      \item El SJM no logra delimitar o deserializar el mensaje recibido conforme al protocolo acordado.
      \item El SJM descarta el mensaje, cierra la conexión y registra en la bitácora local del cliente el contenido truncado y la posición donde falló la lectura.
      \item El SJM muestra la \texttt{GUI\_\allowbreak{}MessageError} indicando que la respuesta del servidor no pudo interpretarse.
      \item Fin del caso de uso.
    \end{cuenum}} \\ \hline
  \textbf{Postcondiciones} & \cuCelda{\begin{cureglas}
      \item \textbf{\mbox{POST-01}} Si el caso termina por el FN, existe una \textit{Sesion} activa asociada al \textit{Usuario}, con su token, su \texttt{fecha\_\allowbreak{}inicio}, la dirección de red de origen y la huella del dispositivo. Las demás sesiones abiertas de la cuenta, si las hay, siguen abiertas (\mbox{D-01}).
      \item \textbf{\mbox{POST-02}} Si el caso termina por el FN, \texttt{intentos\_\allowbreak{}fallidos} del \textit{Usuario} vale cero y \texttt{bloqueada\_\allowbreak{}hasta} es nulo.
      \item \textbf{\mbox{POST-03}} Si el caso termina por el FN, el último acceso de la cuenta es la \texttt{fecha\_\allowbreak{}inicio} de la nueva \textit{Sesion}; no se copia en el \textit{Usuario}.
      \item \textbf{\mbox{POST-04}} Si el caso termina por el FN, el cliente conserva el token de sesión y muestra la \texttt{GUI\_\allowbreak{}MainMenu} con la interfaz en el \texttt{idioma\_\allowbreak{}preferido} del \textit{Usuario}.
      \item \textbf{\mbox{POST-05}} Si el caso termina por el \mbox{FA-08}, \texttt{intentos\_\allowbreak{}fallidos} quedó incrementado, \texttt{fecha\_\allowbreak{}ultimo\_\allowbreak{}intento\_\allowbreak{}fallido} actualizado y, de haberse alcanzado un umbral, \texttt{bloqueada\_\allowbreak{}hasta} contiene el instante en que expira el bloqueo.
      \item \textbf{\mbox{POST-06}} Si el caso termina por cualquier flujo alterno o excepción distinta del FN, del \mbox{FA-17} o del \mbox{FA-18}, no existe ninguna \textit{Sesion} nueva y el Jugador permanece en la \texttt{GUI\_\allowbreak{}Login} o en la pantalla informativa correspondiente.
      \item \textbf{\mbox{POST-07}} En todos los casos que llegan al servidor, exitosos o fallidos, el intento quedó registrado en la \textit{BitacoraAcceso} con su resultado, su fecha y hora y la dirección de red de origen.
    \end{cureglas}} \\ \hline
  \textbf{Reglas de negocio} & \cuCelda{\begin{cureglas}
      \item \textbf{\mbox{RN-01}} Ante credenciales inválidas, el sistema responde siempre con el mismo mensaje genérico y en un tiempo equivalente, sin revelar si el fallo fue el identificador o la contraseña, ni si la cuenta existe.
      \item \textbf{\mbox{RN-02}} Las contraseñas y los códigos del segundo factor no se almacenan, transmiten ni registran en claro. En la base de datos se guarda únicamente su resumen criptográfico con sal, y ninguna bitácora puede contener su valor.
      \item \textbf{\mbox{RN-03}} El bloqueo por intentos fallidos es escalonado: al tercer intento fallido consecutivo la cuenta se bloquea cinco minutos; al sexto, treinta minutos; a partir del noveno, sesenta minutos. El contador se restablece con un acceso exitoso, o bien cuando transcurren veinticuatro horas desde \texttt{fecha\_\allowbreak{}ultimo\_\allowbreak{}intento\_\allowbreak{}fallido} sin nuevos intentos.
      \item \textbf{\mbox{RN-04}} Solo puede iniciar sesión un \textit{Usuario} cuyo \texttt{estado\_\allowbreak{}cuenta} sea \texttt{ACTIVA}. Los valores \texttt{PENDIENTE}, \texttt{SUSPENDIDA} y \texttt{BANEADA} impiden el acceso, cada uno con su propio mensaje. Una cuenta \texttt{ELIMINADA} tiene el \texttt{nickname} y el \texttt{correo} anonimizados, por lo que nunca coincide con el identificador recibido y se trata como inexistente (\mbox{FA-06}).
      \item \textbf{\mbox{RN-05}} Una \textit{Sancion} activa de ámbito \texttt{CUENTA} impide iniciar sesión hasta su \texttt{fecha\_\allowbreak{}fin}. Si la sanción es permanente, \texttt{fecha\_\allowbreak{}fin} es nula y el impedimento no caduca. Una sanción de ámbito \texttt{CHAT} no afecta al inicio de sesión (\mbox{D-05}).
      \item \textbf{\mbox{RN-06}} Una cuenta puede tener varias \textit{Sesion} abiertas a la vez, una por dispositivo. Lo que no puede duplicarse es la participación en una partida, que se controla al entrar a ella y no aquí (\mbox{D-01}).
      \item \textbf{\mbox{RN-07}} El token de sesión tiene una vigencia máxima de veinticuatro horas y se invalida al cerrar sesión, al cambiar las credenciales o al aplicarse una sanción de ámbito de cuenta.
    \end{cureglas}} \\
   & \cuCelda{\begin{cureglas}
      \item \textbf{\mbox{RN-08}} El segundo factor se solicita únicamente cuando el \textit{Usuario} lo tiene activado en su cuenta. El mecanismo concreto de entrega y verificación queda detrás de una interfaz, de modo que cambiarlo no obligue a modificar este flujo.
      \item \textbf{\mbox{RN-09}} El código del segundo factor tiene una vigencia de cinco minutos, admite un máximo de tres intentos de verificación y es de un solo uso. Solo puede existir un código vigente por cuenta.
      \item \textbf{\mbox{RN-10}} El correo de verificación puede reenviarse a lo sumo una vez cada cinco minutos por cuenta.
      \item \textbf{\mbox{RN-11}} El tiempo transcurrido entre el paso 2 y el paso 19 del FN debe ser menor a un segundo, sin contar el tiempo que consuman los servicios externos, acotados por los tiempos límite descritos en \mbox{EX-03} y \mbox{EX-06}.
      \item \textbf{\mbox{RN-12}} Todos los textos visibles de las pantallas de este caso provienen de los catálogos de recursos y no del código, y se almacenan y transmiten en Unicode. Los mensajes de error también se localizan.
      \item \textbf{\mbox{RN-13}} Las validaciones que el cliente realiza sin consultar al servidor no cuentan como intentos fallidos y no afectan \texttt{intentos\_\allowbreak{}fallidos}.
      \item \textbf{\mbox{RN-14}} El \texttt{nickname} y el \texttt{correo} son únicos entre todos los \textit{Usuario}, ya que ambos sirven como identificador de acceso.
    \end{cureglas}} \\
   & \cuCelda{\begin{cureglas}
      \item \textbf{\mbox{RN-15}} Las cuentas de invitado no inician sesión por este caso: no tienen \texttt{correo} ni contraseña y se recuperan con el token guardado en el dispositivo (\mbox{CU-06}).
      \item \textbf{\mbox{RN-16}} Una sanción temporal de cuenta no necesita un proceso que la levante: el \texttt{estado\_\allowbreak{}cuenta} vuelve a \texttt{ACTIVA} en el primer inicio de sesión posterior a su \texttt{fecha\_\allowbreak{}fin} (\mbox{FA-19}). El valor \texttt{SUSPENDIDA} indica que hubo una suspensión, pero la vigencia la decide siempre la \textit{Sancion}.
    \end{cureglas}} \\ \hline
  \textbf{Trazabilidad} & \cuCelda{Restricciones \mbox{CON-01} (C\#), \mbox{CON-02} y \mbox{CON-03} (TCP sin WebSockets), \mbox{CON-04} (SQL Server), \mbox{CON-05} (Unicode), \mbox{CON-07} (respuesta menor a un segundo), \mbox{CON-08} (segundo factor), \mbox{CON-11} (cambio de idioma) y \mbox{CON-12} (comprensible a los ocho años). Concerns \mbox{CRN-05}, \mbox{CRN-21}, \mbox{CRN-22} y \mbox{CRN-26}. Características de calidad \mbox{QA-01}, \mbox{QA-02}, \mbox{QA-03}, \mbox{QA-04}, \mbox{QA-05}, \mbox{QA-06} y \mbox{QA-08}. Preguntas abiertas \mbox{Q-03} y \mbox{Q-04}, resueltas para este caso como se indica en las observaciones.} \\ \hline
  \textbf{Observación} & \cuCelda{\textit{Sobre el segundo factor opcional por cuenta (\mbox{Q-03}).}\par
    El segundo factor se solicita solo si el \textit{Usuario} lo activó, lo que se representa con \texttt{doble\_\allowbreak{}factor\_\allowbreak{}habilitado}. La restricción que exige su existencia se cumple porque el mecanismo está construido y disponible; se deja a criterio de cada jugador porque un niño de ocho años difícilmente administra una aplicación autenticadora, y obligarlo cerraría el acceso a buena parte del público objetivo. El mecanismo de entrega queda detrás del Proveedor del segundo factor, de modo que sustituirlo no toca este caso de uso.} \\ \hline
  \textbf{Observación} & \cuCelda{\textit{Sobre el bloqueo temporal escalonado.}\par
    Se eligió frente al bloqueo permanente porque no introduce al servicio de correo en la ruta de inicio de sesión y porque no deja fuera de su cuenta a un jugador que simplemente olvidó su contraseña. El escalonamiento encarece progresivamente la fuerza bruta sin castigar el error honesto.} \\ \hline
  \textbf{Observación} & \cuCelda{\textit{Sobre el acceso de cuentas no verificadas (\mbox{Q-04}).}\par
    Una cuenta en \texttt{PENDIENTE} no puede iniciar sesión. La alternativa de permitir un acceso con funciones limitadas se descartó porque obligaría a que todos los demás casos de uso comprobaran el estado de la cuenta antes de operar, con el costo de mantenimiento que eso implica.} \\ \hline
  \textbf{Observación} & \cuCelda{\textit{Sobre las sesiones simultáneas (\mbox{D-01}).}\par
    La versión anterior de este caso limitaba cada cuenta a una sola sesión activa y obligaba a cerrar la del otro dispositivo para poder entrar. La restricción existía para que dos clientes no enviaran jugadas en nombre del mismo jugador, pero ese riesgo está en la partida y no en el inicio de sesión. Se trasladó al emparejamiento como regla de una sola participación sin terminar por jugador, y con ella desaparecieron el flujo alterno de sesión ya abierta y el motivo de cierre por desplazamiento.} \\ \hline
  \textbf{Datos que toca} & \cuCelda{\begin{cudatos}
      \item \textbf{\textit{Usuario}:} lee \texttt{nickname}, \texttt{correo}, \texttt{tipo\_\allowbreak{}cuenta}, \texttt{contrasena\_\allowbreak{}hash}, \texttt{contrasena\_\allowbreak{}sal}, \texttt{bloqueada\_\allowbreak{}hasta}, \texttt{estado\_\allowbreak{}cuenta}, \texttt{doble\_\allowbreak{}factor\_\allowbreak{}habilitado}, \texttt{idioma\_\allowbreak{}preferido}, \texttt{id\_\allowbreak{}icono}, \texttt{fecha\_\allowbreak{}configuracion\_\allowbreak{}inicial} \textit{(pasos 6, 7, 8, 9, 11, 13, 17)}
      \item \textbf{\textit{Usuario}:} escribe \texttt{intentos\_\allowbreak{}fallidos}, \texttt{fecha\_\allowbreak{}ultimo\_\allowbreak{}intento\_\allowbreak{}fallido}, \texttt{bloqueada\_\allowbreak{}hasta}, \texttt{fecha\_\allowbreak{}ultimo\_\allowbreak{}envio\_\allowbreak{}verificacion}; \texttt{estado\_\allowbreak{}cuenta} en el \mbox{FA-19} \textit{(pasos 10, \mbox{FA-08}, \mbox{FA-09}, \mbox{FA-19})}
      \item \textbf{\textit{Sancion}:} lee las activas de ámbito \texttt{CUENTA} \textit{(pasos 12, \mbox{FA-10})}
      \item \textbf{\textit{Sesion}:} crea \textit{(paso 14)}
      \item \textbf{\textit{CodigoSegundoFactor}:} crea, marca como consumido o invalidado, incrementa \texttt{intentos} \textit{(pasos \mbox{FA-11}, \mbox{FA-13} a \mbox{FA-16})}
      \item \textbf{\textit{TokenVerificacionCorreo}:} invalida los pendientes y crea uno nuevo \textit{(paso \mbox{FA-09})}
      \item \textbf{\textit{EstadisticaModo}, \textit{Equipamiento}, \textit{Participacion}:} lee, para el mensaje \texttt{LOGIN\_\allowbreak{}OK} \textit{(paso 17)}
      \item \textbf{\textit{BitacoraAcceso}:} inserta \textit{(pasos 15, \mbox{FA-06} a \mbox{FA-10}, \mbox{FA-13}, \mbox{FA-16}, \mbox{FA-19})}
    \end{cudatos}} \\ \hline
\end{fichacu}
\clearpage
